\section{Related Work}

\paragraph{Evidence-grounded long-form generation.}
Retrieval-augmented generation improves the factuality and verifiability of language models by grounding generation in external evidence~\citep{lewis2020retrieval}. Recent work extends this paradigm to web-grounded and citation-grounded long-form generation, where models must synthesize evidence from multiple sources and provide faithful attributions~\citep{liu2023webglm,qin2023webcpm,gao2023enabling}. Scientific and legal domains further stress these abilities, as reliable generation often requires long-document reading, evidence comparison, and well-supported argumentation~\citep{openscholar,mcdonald2022detect,guha2023legalbench,dai2025laiw}. However, most existing pipelines supervise models through final answers, retrieved contexts, or distilled responses, collapsing the intermediate evidence-seeking process into a single input-output mapping. \textsc{RetroGen} instead treats expert-written artifacts as compressed traces of latent evidence-seeking processes and reconstructs explicit tool-use trajectories from them.

\paragraph{Tool-use agents and process supervision.}
A growing line of work studies language models as agents that interleave reasoning with external actions such as search, retrieval, API calls, and environment interaction~\citep{wang2024survey,yao2022react,qin2024toolllm}. While tool use extends models beyond parametric knowledge, training reliable agents remains challenging because high-quality process trajectories are scarce. Process supervision shows that intermediate reasoning steps can improve reliability over outcome-only supervision~\citep{lightman2024let}, and reasoning bootstrapping suggests that models can benefit from learning from their own intermediate rationales~\citep{zelikman2022star,shao2024deepseekmath}. Unlike these settings, evidence-seeking agents require not only plausible reasoning chains but also faithful observations produced by executable tools. \textsc{RetroGen} therefore reconstructs candidate trajectories retrospectively and filters them with evidence-constrained verification.

\paragraph{Self-improvement without stronger teachers.}
Self-improving agents have been explored through verbal feedback, reflection, and iterative refinement~\citep{shinn2023reflexion}. However, many data-generation pipelines still rely on stronger teacher models, human demonstrations, or preference labels to produce supervision at scale. Web-grounded systems such as WebGLM and WebCPM show the value of retrieval-enhanced supervision for long-form QA~\citep{liu2023webglm,qin2023webcpm}, but typically organize supervision around retrieve-then-write outputs rather than reusable multi-step agentic processes. \textsc{RetroGen} differs by using expert-curated artifacts as target anchors: instead of asking a stronger model or human annotator to generate trajectories forward, it infers plausible evidence-seeking trajectories backward from the artifact, verifies them automatically, and trains the backbone on the resulting retrospective process supervision.